\PassOptionsToPackage{dvipsnames,svgnames,table}{xcolor}
\documentclass[11pt]{article}
\usepackage[a4paper,margin=22mm]{geometry}
\usepackage[no-math]{fontspec}
\setmainfont{DejaVu Serif}
\usepackage{amsmath,amssymb,amsthm,mathtools,graphicx,xcolor,hyperref,xurl,xspace,ifthen,enumerate}
\usepackage[shortlabels]{enumitem}
\setlength{\emergencystretch}{4em}
\AtBeginDocument{\special{pdf:mapline pzdr Dingbats <uzdr.pfb}}
\SetMathAlphabet{\mathbf}{normal}{OT1}{cmr}{bx}{n}
\SetMathAlphabet{\mathbf}{bold}{OT1}{cmr}{bx}{n}
\usepackage[nameinlink]{cleveref}
\newtheorem{theo}{Theorem}[section]
\newtheorem{lemm}[theo]{Lemma}
\newtheorem{thm}[theo]{Theorem}
\newtheorem{proposition}[theo]{Proposition}
\newtheorem{corollary}[theo]{Corollary}
\newtheorem{definition}[theo]{Definition}
\newtheorem{theorem}[theo]{Theorem}
\newtheorem*{remas*}{Remarks}
\newtheorem{remas}[theo]{Remarks}
\newtheorem*{exam*}{Example}
\newtheorem{ques}[theo]{Question}
\providecommand{\og}{«\nobreakspace}
\providecommand{\fg}{\nobreakspace»}
\newtheorem{lem}[theo]{Lemma}
\newtheorem{lemma}[theo]{Lemma}
\newtheorem{prop}[theo]{Proposition}
\newtheorem{coro}[theo]{Corollary}
\newtheorem{corol}[theo]{Corollary}
\newtheorem*{conj*}{Conjecture}
\newtheorem{conj}[theo]{Conjecture}
\newtheorem{conjecture}[theo]{Conjecture}
\newtheorem{Proposition}[theo]{Proposition}
\newtheorem{theoreme}[theo]{Theorem}
\theoremstyle{definition}
\newtheorem{defi}[theo]{Definition}
\newtheorem{exam}[theo]{Example}
\newtheorem{exem}[theo]{Example}
\newtheorem{remark}[theo]{Remark}
\newtheorem{example}[theo]{Example}
\newtheorem{exams}[theo]{Examples}
\newtheorem{rema}[theo]{Remark}
\newtheorem{nota}[theo]{Notation}
\newtheorem{notation}[theo]{Notation}
\newtheorem*{notas}{Notations}
\newtheorem{result}[theo]{Result}
\newtheorem{question}[theo]{Question}
\newtheorem{prob}[theo]{Problem}
\newtheorem*{rema*}{Remark}
\newtheorem*{defi*}{Definition}
\newtheorem*{theo*}{Theorem}
\newtheorem*{lemm*}{Lemma}
\newtheorem*{prop*}{Proposition}
\newcommand{\enoncename}{}
\newtheorem{corpusenonce}[theo]{\enoncename}
\newenvironment{enonce}[2][]{\renewcommand{\enoncename}{#2}\begin{corpusenonce}}{\end{corpusenonce}}
\newenvironment{enonce*}[2][]{\par\medskip\noindent\textbf{#2.}\itshape}{\par\medskip}
\newenvironment{proofc}[1][\proofname]{\begin{proof}[#1]}{\end{proof}}
\newcommand{\Mc}{\mathcal{M}}
\newcommand{\pointir}{.\ ---\ }
\newcommand{\equalenv}[2]{\expandafter\let\csname#1\expandafter\endcsname\csname#2\endcsname\expandafter\let\csname end#1\expandafter\endcsname\csname end#2\endcsname}
\providecommand{\diagbox}[3][]{\begin{tikzpicture}[x=1em,y=1em,baseline=(current bounding box.center)]\useasboundingbox(0,0) rectangle(3,1.4);\draw(0,1.4)--(3,0);\node at(.5,.3){#2};\node at(2.5,1.1){#3};\end{tikzpicture}}
\providecommand{\eg}{e.g.\ }
\providecommand{\ie}{i.e.\ }
\providecommand{\cf}{cf.\ }
\providecommand{\longthanks}[1]{\par\smallskip\noindent #1\par\smallskip}
\providecommand{\qbinom}[2]{\genfrac{[}{]}{0pt}{}{#1}{#2}}

\providecommand{\mathof}[1]{\mathrm{#1}}

\numberwithin{equation}{section}


\def\endappendix{}
\PassOptionsToPackage{svgnames,dvipsnames}{xcolor}

\usepackage{amsfonts}
\usepackage{amsmath} 
\usepackage{amssymb}
\usepackage{amsthm}
\usepackage{graphicx}
\usepackage{tikz}
\usetikzlibrary{calc}
\usetikzlibrary{intersections}
\usetikzlibrary{arrows.meta}
\usepackage[nameinlink]{cleveref}
\usepackage{booktabs}
\usepackage[title]{appendix}

\Crefname{theo}{Theorem}{Theorems}
\Crefname{coro}{Corollary}{Corollaries}
\Crefname{prop}{Proposition}{Propositions}
\Crefname{lemm}{Lemma}{Lemmas}
\Crefname{defi}{Definition}{Definitions}
\Crefname{rema}{Remark}{Remarks}

\DeclareMathOperator{\rank}{rank}
\DeclareMathOperator{\outdeg}{deg^+}
\DeclareMathOperator{\indeg}{deg^-}
\DeclareMathOperator{\diag}{diag}
\newcommand{\diff}{\mathrm{d}}
\newcommand{\undirected}[1]{\overline{#1}}

\colorlet{colbg}{white}
\colorlet{colfg}{black}
\colorlet{colgraphv}{colfg!75!colbg}
\colorlet{colgraphe}{colfg!55!colbg}
\colorlet{colG}{DarkSeaGreen}
\definecolor{colR}{HTML}{CC6677}
\definecolor{colO}{HTML}{DDCC77}
\definecolor{colB}{HTML}{6699CC}
\colorlet{colY}{Gold!90!black}
\colorlet{colGray}{white!60!black}

\colorlet{col1}{colR}
\colorlet{col2}{colG}
\colorlet{col3}{colGray}
\colorlet{col4}{colO}
\colorlet{col5}{colB}
\colorlet{colr1}{colR}
\colorlet{colr2}{colB}
\colorlet{colr3}{colO}
\colorlet{colr4}{colG}
\colorlet{colr5}{colGray}
\colorlet{colr6}{Tan}
\colorlet{colr7}{Plum}
\colorlet{colr8}{LightSkyBlue}
\colorlet{colr9}{LightGreen}
\colorlet{cola}{colfg!75!colbg}

\tikzstyle{vertex}=[fill=colgraphv,circle,inner sep=0pt, minimum size=4pt]
\tikzstyle{edge}=[line width=1.5pt,colgraphe]

\tikzstyle{axes}=[draw=cola,-{Latex[round,width=3pt]}]
\tikzstyle{aline}=[draw=cola]
\tikzstyle{bline}=[draw=cola!10!colbg]
\tikzstyle{alabelsty}=[cola,font=\scriptsize]
\tikzstyle{labelsty}=[font=\scriptsize]
\tikzstyle{slabelsty}=[font=\tiny]


\newcommand{\midarrow}{\tikz \draw[-triangle 90] (0,0) -- +(.1,0);}

\newcommand{\dedge}[2]{(#1,#2)}

\newcommand{\bnr}[2]{\tikz[baseline=3pt]{\fill[#2] (0,0)rectangle(#1,0.45);\node[anchor=west] at (0,0.225) {#1};}}
\newcommand{\bnrs}[3]{\tikz[baseline=3pt]{\fill[#3] (0,0)rectangle(#1*#2,0.45);\node[anchor=west] at (0,0.225) {#1};}}
\newcommand{\bnrsm}[3]{\tikz[baseline=3pt]{\fill[#3] (0,0)rectangle(#1*#2-#2,0.45);\node[anchor=west] at (0,0.225) {#1};}}

\renewcommand{\theenumi}{(\roman{enumi})}\renewcommand{\labelenumi}{\theenumi}














\begin{document}
\section*{\raggedright 1006: Spherical generic complex realization counts equal Euclidean counts after coning in one higher dimension}
\noindent\textbf{Author:} Dewar, Sean; Grasegger, Georg\par\noindent\textbf{Source:} The number of realisations of a rigid graph in Euclidean and spherical geometries. Algebraic Combinatorics 7 (2024), publisher version, DOI 10.5802/alco.390. \url{https://alco.centre-mersenne.org/articles/10.5802/alco.390/}
\par\noindent\textbf{License:} CC BY 4.0, \url{https://creativecommons.org/licenses/by/4.0/}. The original publisher PDF cover has the explicit license notice. See the saved source/license records.
\par\noindent\textbf{Editorial material note (not author text):} The selected problem is original Theorem1.2 only: the exact equality between generic complex spherical d-realization counts of G and Euclidean(d+1)-realization counts of its cone G*o, with all original definitions/non-rigid/small-complete conventions. Complete source Sections2--5 through1039 retain complex rigidity/normalized varieties, generic fiber counts, spherical setup, full projective-scaling lemma and dense-open radial map argument; complete author AppendicesA/B1761--2043 retain the generic-degree and Euclidean normalized-fiber prototype proofs explicitly cited by the spherical adaptation. Later plane/sphere inequality, repeated-coning and numerical examples are unselected. All necessary proof text/formulas are editable native TeX; incidental omitted figures are retained only as their original printed references. Literal source-context note: Y\_\{G,d\} uses points in C\textasciicircum{}\{d+1\}, while selected-support Lemma5.1 at native936 writes p\_v-prime=(r\_v p\_v,0) in X\_\{G*o,d+1\}. The complete forward/reverse human construction at967--989 and main radial map1021--1026 use q\_v-prime=r\_v q\_v and p\_v-prime=r\_v p\_v without that appended zero. The literal unfinished 'since' at962 and U ambient wording at1009 are also retained. Root read the complete definitions, support and primary proof and permitted this bounded source-context disposition. All literal source formulas and wording remain unchanged; no editor coordinate removal, ambient-space correction or added argument is supplied. All proof text and formulas below are editable author TeX. Mathematical wording, language and proof order are retained. Only the article class, title matter and bibliography presentation are adapted for independent compilation; no proof has been generated or repaired.
\par\noindent\textbf{Editorial difficulty:} H3: The author proves projective scaling invariance of normalized coned fibers, identifies exact generic fiber cardinalities for Euclidean and spherical frameworks, and intersects dense open loci through a dominant radial-scaling morphism. This establishes the exact cone/sphere counting identity in every dimension, beyond qualitative rigidity of cones.
\medskip\hrule\medskip
\noindent\textbf{Author statement, complete proof and local supporting text follow in source order.}
\setcounter{section}{1}\setcounter{theo}{1}
\begin{theorem}\label{t:conecount}
	Let $d$ be a positive integer and let $G*o$ be a coning of a graph $G$.
	Then $c^*_{d}(G) = c_{d+1}(G*o)$.
\end{theorem}

\section{Preliminary results for rigidity theory}\label{sec:prelim}

In this section we cover the necessary background results for discussing rigidity in geometries based over the real numbers.
Many of these concepts are extended to geometries based over the complex numbers in \Cref{sec:count,sec:countsphere}.

\subsection{Euclidean space rigidity}

A \emph{realisation} of a (finite simple) graph $G=(V,E)$ in $\mathbb{R}^d$ is a map $p:V \rightarrow \mathbb{R}^d$,
and the linear space of all realisations of a graph is denoted by $(\mathbb{R}^d)^V$.
A realisation $p$ is said to be \emph{generic} if the set of coordinates of $p$ forms an algebraic independent set of $d|V|$ elements.\footnote{The requirement that the algebraic independent set has $d|V|$ elements is stated to avoid the possibility of identical coordinates.}
A graph-realisation pair $(G,p)$ is said to be a \emph{$d$-dimensional framework}.
Two frameworks $(G,p)$ and $(G,q)$ in $\mathbb{R}^d$ are said to be \emph{equivalent} if (given $\|\cdot\|$ is the standard Euclidean norm) $\|p_v - p_w\| = \|q_v-q_w\|$ holds for every edge $vw \in E$.
A framework $(G,p)$ in $\mathbb{R}^d$ is now said to be \emph{rigid} if the following holds for some $\varepsilon >0$:
if $(G,q)$ is an equivalent framework in $\mathbb{R}^d$ such that $\|p_v-q_w\|<\varepsilon$ for all $v \in V$,
then there exists an isometry $f: \mathbb{R}^d \rightarrow \mathbb{R}^d$ such that $q_v = f(p_v)$ for all $v \in V$.

Determining whether a framework is rigid is NP-Hard for $d \geq 2$~\cite{Abbot}.
To combat this, we construct the \emph{rigidity matrix} $R(G,p)$ of a given framework $(G,p)$.
This is the $|E| \times d|V|$ matrix with the row labelled $vw \in E$ given by
\begin{align*}
	[ 0 \quad \cdots  \quad 0  \quad \overbrace{p_v-p_w}^{v} \quad 0  \quad \cdots  \quad 0  \quad \overbrace{p_w-p_v}^{w} \quad 0  \quad \cdots  \quad 0].
\end{align*}
A framework $(G,p)$ is said to be \emph{regular} if its rigidity matrix has maximal rank,
\ie for any $d$-dimensional realisation $q$ of $G$ we have $\rank R(G,p) \geq \rank R(G,q)$.
Any framework $(G,p)$ with a surjective rigidity matrix (\ie $\rank (G,p)=|E|$) is automatically regular:
in such a case we say that the framework $(G,p)$ is \emph{independent}.
Any generic framework must also be regular, since the non-regular frameworks form an algebraic set defined by rational-coefficient polynomials.


\begin{theorem}[\cite{AsimowRothI}]\label{t:asimowroth}
	Let $(G,p)$ be a regular framework in $\mathbb{R}^d$.
	Then the following properties are equivalent.
	\begin{enumerate}
		\item \label{t:asimowroth1} $(G,p)$ is rigid.
		\item \label{t:asimowroth2} Either $G$ has at least $d$ vertices and $\rank R(G,p) = d|V| - \binom{d+1}{2}$, or $G$ is a complete graph.
	\end{enumerate}
\end{theorem}

Since every generic framework is regular,
\Cref{t:asimowroth} informs us that both rigidity and independence are generic properties.
This motivates the following definitions.

\begin{definition}
	A graph $G$ is said to be \emph{$d$-rigid} (respectively, $d$-independent) if there exists a $d$-dimensional rigid (respectively, independent) generic framework $(G,p)$.
	If $G$ is rigid but $G-e$ is not for each edge $e \in E$,
	then $G$ is said to be \emph{minimally $d$-rigid}.
\end{definition}

One method that can be used to construct rigid/independent graphs in higher dimensions is via coning.
Specifically:
given a graph $G=(V,E)$,
we define the \emph{cone of $G$ (by the new vertex $o$)} to be the graph $G*o := (V*o, E*o)$ where $o \notin V$, $V*o := V \cup \{o\}$ and 
\begin{align*}
	E*o := E \cup \{ o v : v \in V\}.
\end{align*}
We can now use the next result to move between dimensions whilst preserving rigidity and independence.

\begin{theorem}[{\cite[Theorem 5]{coning}}]\label{t:cone}
	A graph $G=(V,E)$ is $d$-rigid \textup{(}respectively, $d$-independent\textup{)} if and only if $G*o$ is $(d+1)$-rigid \textup{(}respectively, $(d+1)$-independent\textup{)}<.
\end{theorem}
 




\subsection{Spherical space rigidity}

The \emph{$d$-dimensional (unit) sphere} is the set
\begin{align*}
	\mathbb{S}^d := \left\{ x \in \mathbb{R}^{d+1} : \|x\|^2=1 \right\}.
\end{align*}
Similar to the Euclidean case, we define a \emph{$d$-dimensional spherical realisation} to be a pair $(G,p)$ of a graph $G=(V,E)$ and a \emph{$d$-dimensional spherical realisation} \mbox{$p: V \rightarrow \mathbb{S}^d$}.
The notions of equivalence and rigidity can be analogously defined for spherical frameworks by restricting all realisations to be spherical realisations.

An important observation is that any $d$-dimensional spherical realisation of a graph $G$ can (for rigidity purposes at least) be considered to be a $(d+1)$-dimensional realisation of the cone of $G$.
To be more specific:
given a $d$-dimensional spherical framework $(G,p)$,
we define the $(d+1)$-dimensional framework $(G*o,p')$ by setting $p'_v=(p_v,1)$ for all $v \in V$ and $p'_o = \mathbf{0}$.
From this construction it is easy to see that the framework $(G*o,p')$ is rigid if and only if $(G,p)$ is rigid.
Slightly less obviously,
this remains true if we scale each $p_v$ by some positive scalar.

\begin{proposition}\label{p:slide}
	For a graph $G=(V,E)$,
	let $p,q$ be $d$-dimensional realisations of the cone $G*o$ where $p_o=q_o = \mathbf{0}$.
	Further suppose that for each $v \in V$,
	there exists a scalar $r_v >0$ such that $q_v = r_v p_v$.
	Then $(G*o,p)$ is rigid if and only if $(G*o,q)$ is rigid,
	and $\rank R(G*o,p) = \rank R(G*o,q)$.
\end{proposition}

\begin{proof}
	That $\rank R(G*o,p) = \rank R(G*o,q)$ follows from the observation that the rank of a rigidity matrix is projectively invariant;
	see~\cite[Theorem 1]{Izm09} for more details.
	Now choose any $d$-dimensional framework $(G*o,\tilde{p})$ equivalent to $(G*o,p)$ and define $\tilde{q}$ to be the realisation of $G*o$ with $\tilde{q}_v = r_v \tilde{p}_v$ for each $v \in V$.
	For every edge $vw$ of $G$ we have
	\begin{align*}
		\|\tilde{p}_v\|^2 &= \|\tilde{p}_v - \tilde{p}_o\|^2 = \|p_v - p_o\|^2 = \|p_v\|^2,\\
		\|\tilde{p}_w\|^2 &= \|\tilde{p}_w - \tilde{p}_o\|^2 = \|p_w - p_o\|^2 = \|p_w\|^2,\\
		\tilde{p}_v \cdot \tilde{p}_v &=  \|\tilde{p}_v\|^2 + \|\tilde{p}_w\|^2 - \| \tilde{p}_v - \tilde{p}_w\|^2= \|p_v\|^2 + \|p_w\|^2 - \| p_v - p_w\|^2 = p_v \cdot p_w.
	\end{align*}
	(Remember that $ov , ow \in E*o$ and $\tilde{p}_o = p_o = \mathbf{0}$.)
	Using the above equalities,
	we see that
	\begin{align*}
		\|\tilde{q}_v-\tilde{q}_w\|^2 &= \|r_v\tilde{p}_v-r_w\tilde{p}_w\|^2 \\
		&= \|r_v\tilde{p}_v\|^2 + \|r_w\tilde{p}_w\|^2 - r_v r_w\tilde{p}_v \cdot \tilde{p}_w \\
		&= \|r_vp_v\|^2 + \|r_wp_w\|^2 - r_v r_wp_v \cdot p_w  \\
		&= \|r_vp_v-r_wp_w\|^2 \\
		&= \|q_v-q_w\|^2,
	\end{align*}
	and hence $(G*o,q)$ and $(G*o,\tilde{q})$ are equivalent.
	From this it follows that $(G*o,p)$ is rigid if and only if $(G*o,q)$ is rigid.
\end{proof}

From a combination of \Cref{t:asimowroth}, \Cref{t:cone} and \Cref{p:slide},
we can now see that a graph is rigid in $d$-dimensional Euclidean space if and only if it is rigid in $d$-dimensional spherical space.

\begin{theorem}\label{t:sphereclassic}
	For any graph $G$, the following properties are equivalent.
	\begin{enumerate}
		\item $G$ is $d$-rigid.
		\item Almost every (\ie with Lebesgue measure zero complement) $d$-dimensional spherical framework $(G,p)$ is rigid.
	\end{enumerate}
\end{theorem}

\begin{remark}
	It was first observed by Pogorelov~\cite[Chapter V]{Pog73} that the linear space of infinitesimal motions of a $d$-dimensional spherical framework is isomorphic to the linear space of infinitesimal motions of the $d$-dimensional framework obtained by a central projection to Euclidean space.
	An alternative proof of \Cref{t:sphereclassic} now stems from the spherical analogue to \Cref{t:asimowroth},
	which can be proven in an almost identical way.
\end{remark}


\section{Counting complex realisations}\label{sec:count}

Our aim in this section is to prove that the definition of the $d$-realisation number for a graph alluded to in the introduction can be stated in a rigorous manner (see \Cref{def:count}).

\subsection{Complex rigidity map}

We recall that an algebraic set is a subset $A \subset \mathbb{C}^n$ of common zeroes of an ideal $I \subset \mathbb{C}[X_1,\ldots,X_n]$,
and a variety is an irreducible algebraic set.
If an algebraic set forms a smooth submanifold of $\mathbb{C}^n$ then it is said to be smooth.
We shall refer to a subset $U$ of an algebraic set $A$ as a Zariski closed/open/dense subset if $U$ is a closed/open/dense subset of $A$ with respect to the Zariski topology.

For every point  $x=(x_1,\ldots,x_d) \in \mathbb{C}^d$,
we define $[x]_i := x_i$ for each $i =1,\ldots,d$.
We now also consider realisations in $\mathbb{C}^d$ by extending $(\mathbb{R}^d)^V$ to the set $(\mathbb{C}^d)^V$.
Any realisation in $(\mathbb{R}^d)^V$ is said to be a \emph{real realisation}.
We extend the square of the Euclidean norm to $\mathbb{C}^d$ by defining $\|x\|^2 := \sum_{i=1}^d [x]_i^2$.
We note that this is a quadratic form and not the square of the standard complex norm,
and so can take complex values.
Later we shall also require the bilinear map $x \cdot y := \sum_{i =1}^d [x]_i [y]_i$,
where we notice that $x \cdot x = \| x\|^2$.

For any graph $G=(V,E)$ we define the \emph{complex rigidity map} to be the multivariable map
\begin{align*}
	f_{G,d} : (\mathbb{C}^{d})^V \rightarrow \mathbb{C}^E, ~ p \mapsto \left( \frac{1}{2}\|p_v-p_w\|^2 \right)_{vw \in E}.
\end{align*}
We denote the Zariski closure of the image of $f_{G,d}$ by $\ell_d(G)$.
Since the domain of $f_{G,d}$ is irreducible,
$\ell_d(G)$ is a variety.
Note that two realisations $p,q$ of $G$ in $\mathbb{R}^d$ are equivalent if and only if $f_{G,d}(p)=f_{G,d}(q)$.
Given $O(d, \mathbb{C})$ is the group of $d \times d$ complex-valued matrices $M$ where $M^T M = M M^T = I$,
we define two realisations $p,q \in (\mathbb{C}^{d})^V$ to be \emph{congruent} (denoted by $p \sim q$) if and only if there exists $A \in O(d, \mathbb{C})$ and $x \in \mathbb{C}^{d}$ so that $p_v = Aq_v + x$ for all $v \in V$.
If the set of vertices of $(G,p)$ affinely span $\mathbb{C}^d$,
we have the following equivalent statement:
two realisations $p,q$ are congruent if and only if $f_{K_V,d}(p) = f_{K_V,d}(q)$ (see~\cite[Section 10]{Gortler2014} for more details).
For all $p \in (\mathbb{C}^{d})^V$,
we define 
\begin{align*}
	C_d(G,p) := f^{-1}_{G,d} (f_{G,d}(p))/_\sim
\end{align*}
to be the \emph{realisation space of $(G,p)$}.

Our new definitions might look a little strange at first.
For example,
the elements of $O(d, \mathbb{C})$ are not isometries of $\mathbb{C}^d$ (\ie they are not unitary matrices),
although $\|Ax\|^2 = \|x\|^2$ for all $x \in \mathbb{C}^d$.
However, our previous definitions of independence and rigidity can be encoded in our new language of morphisms between complex spaces.
We first recall that a morphism $f : X \rightarrow Y$ between algebraic sets $X \subset \mathbb{C}^m$ and $Y \subset \mathbb{C}^n$ (\ie~the restriction of a polynomial map $\mathbb{C}^m \rightarrow \mathbb{C}^n$) is \emph{dominant} if $Y \setminus f(X)$ is contained in a Zariski closed proper subset of $Y$;
see \Cref{sec:dom} for more details regarding dominant morphisms.


\begin{lemma}\label{l:domind}
	Let $G=(V,E)$ be any graph.
	Then the following are equivalent:
	\begin{enumerate}
		\item \label{l:domind1} $G$ is $d$-independent.
		\item \label{l:domind2} The map $f_{G,d}$ is dominant.
		\item \label{l:domind3} $\ell_d(G) = \mathbb{C}^E$.
	\end{enumerate}
\end{lemma} 

\begin{proof}
	By the definition of a dominant map, \ref{l:domind2} and \ref{l:domind3} are equivalent.
	The map $f_{G,d}$ is dominant if and only if $\rank \diff f_{G,d}(p) = |E|$ for some $p \in (\mathbb{C}^{d})^V$ (see \Cref{borel91}).
	If $G$ is $d$-independent then there exists a $p \in (\mathbb{R}^{d})^V$ such that $\rank \diff f_{G,d}(p) = \rank R(G,p) = |E|$,
	hence \ref{l:domind1} implies \ref{l:domind2}.
	Suppose that $G$ is not $d$-independent;
	i.\,e., for each $p \in (\mathbb{R}^{d})^V$ we have $\rank \diff f_{G,d}(p) < |E|$.
	Then the set $X := \{ p \in (\mathbb{C}^{d})^V: \rank \diff f_{G,d}(p) < |E| \}$ is a Zariski closed subset of $(\mathbb{C}^{d})^V$ that contains $(\mathbb{R}^{d})^V$.
	Hence, $X = (\mathbb{C}^{d})^V$,
	and \ref{l:domind2} implies \ref{l:domind1} as required.
\end{proof}

Let $G=(V,E)$ be a graph with at least $d+1$ vertices and fix a sequence of $d$ vertices $v_1,\ldots, v_d$.
We now define the algebraic set
\begin{align}\label{eq:xset}
	X_{G,d} := \left\{ p \in (\mathbb{C}^{d})^V : [p_{v_k}]_j=0 \text{ if } j \geq k  \right\}\,.
\end{align}
(At certain points in the paper we use vertices $w_1,\ldots,w_d$ to define $X_{G,d}$,
which can be done by replacing each $v_i$ with $w_i$ in the above definition.)
We further note that $X_{G,d}$ has dimension $d|V| - \binom{d+1}{2}$.
Since $X_{G,d}$ is defined by a set of linear equations,
it is irreducible.
With this, we define the morphism
\begin{align*}
	\tilde{f}_{G,d} : X_{G,d} \rightarrow \ell_d(G), ~ p \mapsto f_{G,d}(p),
\end{align*}
i.\,e., the restriction of $f_{G,d}$ to the domain $X_{G,d}$ and the codomain $\ell_d(G)$.
As can be seen by the following lemma, the map $\tilde{f}_{G,d}$ allows us to more easily define the cardinality of the set $C_d(G,p)$ for most realisations $p$.

\begin{lemma}\label{l:xgd}
	Let $G=(V,E)$ be a graph with $|V| \geq d+1$,
	and fix a sequence of $d$ vertices $v_1,\ldots, v_d$.
	Then the image of $\tilde{f}_{G,d}$ is Zariski dense in the image of $f_{G,d}$ (and so $\tilde{f}_{G,d}$ is dominant),
	and
	\begin{align*}
		\left|\tilde{f}_{G,d}^{-1}\left(f_{G,d}(p) \right) \right| = 2^{d} |C_d(G,p)|
	\end{align*}
	for almost all\footnote{Here we say that a property $P$ holds \emph{for almost all points in $\mathbb{C}^n$} if there exists a Zariski open subset $U \subset \mathbb{C}^n$ of points where property $P$ holds, in which case we say that any point where property $P$ holds is a \emph{general point (with respect to $P$)}.} $p \in (\mathbb{C}^d)^V$.
\end{lemma}

See \Cref{sec:xgd} for a detailed proof of \Cref{l:xgd}.




\subsection{Defining the \texorpdfstring{$d$-realisation}{d-realisation} number}

Before stating our next result, we require the following technical result regarding dominant morphisms.
The proof of this can be seen in \Cref{sec:dom}.

\begin{theorem}\label{t:deg}
	Let $X \subset \mathbb{C}^m$ and $Y \subset \mathbb{C}^n$ be varieties and $f:X \rightarrow Y$ be a dominant morphism.
	Then the following are equivalent:
	\begin{enumerate}
		\item \label{t:deg1} $\dim X = \dim Y$.
		\item \label{t:deg2} $f$ is \emph{generically finite},
		\ie for a general $y \in Y$,
		the fibre $f^{-1}(y)$ is a finite set.
		\item \label{t:deg3} There exists a $k \in \mathbb{N}$ and a non-empty Zariski open subset $U \subset X$ where $|f^{-1}(f(x))| = k$ for all $x \in U$.
		Furthermore,
		if $x \in U$ and $x' \in f^{-1}(f(x))$,
		then $\rank \diff f(x') = \dim Y$.
	\end{enumerate}
\end{theorem}

\begin{proposition}\label{p:dom}
	Let $G=(V,E)$ be a graph with $|V| \geq d+1$.
	Then the following are equivalent:
	\begin{enumerate}
		\item\label{p:dom1} $G$ is $d$-rigid.
		\item\label{p:dom2} The map $\tilde{f}_{G,d}$ is generically finite.
		\item\label{p:dom3} There exists an $n \in \mathbb{N}$ and a non-empty Zariski open subset $U \subset (\mathbb{C}^d)^V$ where $ |C_d(G,p)|=n$ for all $p \in U$.
		Furthermore,
		if $p \in U$ and $q$ is an equivalent $d$-dimensional realisation of $G$,
		then $\rank \diff f_{G,d}(q) =  d|V| - \binom{d+1}{2}$.
	\end{enumerate}
\end{proposition} 

\begin{proof}
	By \Cref{l:xgd},
	the map $\tilde{f}_{G,d}$ is dominant.	
	Since $\dim \ell_d(G) = \rank \diff f_{G,d}(p)$ for any general realisation $p\in (\mathbb{C}^d)^V$,
	it follows from \Cref{t:asimowroth} that $G$ is $d$-rigid if and only if 
	\begin{align*}
		\dim \ell_d(G) = d|V|- \binom{d+1}{2} = \dim X_{G,d}.
	\end{align*}
	The result now holds by applying \Cref{t:deg} to the map $\tilde{f}_{G,d}$.
\end{proof}

Using \Cref{p:dom} we can now make the following well-defined definition of the $d$-realisation number for any graph.

\begin{definition}\label{def:count}
	The \emph{$d$-realisation number} of a graph $G=(V,E)$ is an element of $\mathbb{N} \cup \{\infty\}$ given by
	\begin{align*}
		c_d(G) :=		
		\begin{cases}
			|C_d(G,p)| \text{ for a general $p \in (\mathbb{C}^d)^V$} &\text{if } |V| \geq d+1, \\
			 1 &\text{if } |V| \leq d \text{ and $G$ is complete}, \\
			 \infty &\text{if } |V| \leq d \text{ and $G$ is not complete}.
		\end{cases}
	\end{align*}
\end{definition}

It follows from \Cref{p:dom} that a graph $G$ is $d$-rigid if and only if $c_d(G) < \infty$.

We can relate the $d$-realisation number back to real realisations so long as we allow for equivalent complex realisations when counting.
We first require the following technical lemma.

\begin{lemma}\label{l:realzar}
	If $U$ is a non-empty Zariski open subset of $\mathbb{C}^n$, then $U \cap \mathbb{R}^n$ is a non-empty Zariski open subset of $\mathbb{R}^n$.
\end{lemma}

\begin{proof}
	For a field $\mathbb{F} \in \{ \mathbb{R}, \mathbb{C}\}$,
	we fix $\mathbb{F}[X^n]$ to be the set of all $n$-variable polynomials over $\mathbb{F}$.
	Define the ideal 
	\begin{align*}
		F := \{ f \in \mathbb{C}[X^n] : f(p) = 0 \text{ for all } p \notin U\}.
	\end{align*}
	Since $U$ is a Zariski open set,
	the zero set
	\begin{align*}
		Z(F) := \{p \in \mathbb{C}^n: f(p) = 0 \text{ for all } f \in F\}
	\end{align*}
	is a Zariski closed set and the complement of $U$.
	As $U$ is non-empty we immediately have that $F \neq \{0\}$.
	For each $f \in F$,
	there exist two unique polynomials \mbox{$f_{\mathcal{R}},f_{\mathcal{I}} \in \mathbb{R}[X^n]$} such that $f = f_{\mathcal{R}}+ i f_{\mathcal{I}}$:
	to see this, note that there exists a finite subset~$A \subset \mathbb{Z}^n_{\geq 0}$ and real values $a_x,b_x $ for each $x \in A$ such that $f(p) = \sum_{x \in A} (a_x + i b_x) p^x$ (here~$p^x = p_1^{x_1} \cdots p_n^{x_n}$ for each $x=(x_1,\ldots,x_n)$),
	and so $f_{\mathcal{R}}(p) = \sum_{x \in A} a_x p^x$ and $f_{\mathcal{I}}(p) = \sum_{x \in A} b_x p^x$.
	For each $p \in \mathbb{R}^n$ we note that $f(p) = 0$ if and only if $f_{\mathcal{R}}(p) = f_{\mathcal{I}}(p) = 0$.
	Fix $F_\mathcal{R} := \{ f_{\mathcal{R}} :f \in F\}$ and~$F_\mathcal{I} := \{ f_{\mathcal{I}} :f \in F\}$,
	and for every set of (real or complex) polynomials $S \subset \mathbb{C}[X^n]$ fix $Z_{\mathbb{R}}(S):= \{p \in \mathbb{R}^n : f(p) = 0 \text{ for all } f \in S \}$.
	Then $Z_{\mathbb{R}}(F) = Z_{\mathbb{R}}(F_{\mathcal{R}}) \cap Z_{\mathbb{R}}(F_{\mathcal{I}})$.
	Hence, the set $Z_{\mathbb{R}}(F)$ is a Zariski closed subset of~$\mathbb{R}^n$ since $Z_{\mathbb{R}}(F_{\mathcal{R}})$ and $Z_{\mathbb{R}}(F_{\mathcal{I}})$ are Zariski closed subsets of~$\mathbb{R}^n$.
	Furthermore,
	the set $Z_{\mathbb{R}}(F)$ is a proper Zariski closed subset of $\mathbb{R}^n$:
	this follows from the observing the equivalence
	\begin{align*}
		Z_{\mathbb{R}}(F) = \mathbb{R}^n \quad \Leftrightarrow \quad F_{\mathcal{R}} = F_{\mathcal{I}} = \{0\}  \quad \Leftrightarrow \quad F= \{0\},
	\end{align*}
	with the latter property contradicting our previous observation that $F \neq \{0\}$.
	The result now holds as $U \cap \mathbb{R}^n = \mathbb{R}^n \setminus Z_{\mathbb{R}}(F)$.
\end{proof}

\begin{corollary}\label{cor:realstuff}
	For each graph $G=(V,E)$,
	there exists a non-empty Zariski open subset $U_\mathbb{R} \subset (\mathbb{R}^d)^V$ of real $d$-dimensional realisations $p$ where $|C_d(G,p)| = c_d(G)$.
\end{corollary}

\begin{proof}
	The result is trivial if $G$ is complete and $|V| \leq d$.	
	If $G$ is not $d$-rigid then $c_d(G) = \infty$ and the result holds by \Cref{t:asimowroth}.
	If $G$ is $d$-rigid with $|V| \geq d+1$ then we fix $U \subset (\mathbb{C}^d)^V$ to be the set defined in \Cref{p:dom} and define $U_{\mathbb{R}} := U \cap (\mathbb{R}^d)^V$.
	The set $U_{\mathbb{R}}$ is now a Zariski open subset of $(\mathbb{R}^d)^V$ by \Cref{l:realzar}.
\end{proof}

It is important to note that \Cref{cor:realstuff} does not guarantee that every other equivalent realisation in $C_d(G,p)$ is also real.
In fact we note that this is often not the case.
For example, any minimally $d$-rigid graph with a vertex of degree $d$ must have general $d$-dimensional realisations with complex equivalent realisations:
this realisation can be chosen in a similar way to the right realisation pictured in original Figure~1.
Failing this,
it is then natural to ask whether each graph $G$ has at least one real generic $d$-dimensional realisation $p$ where every other equivalent realisation in $C_d(G,p)$ is also real.
It was proven in~\cite{JacksonOwen} that this too is false when $d=2$.
It does, however, remain open whether the statement is false when $d>2$.
For more on the topic of equivalent real realisations for a graph,
see~\cite{Bartzos2021}.


When $d=1$, it is relatively easy to compute the realisation number of a graph.
We first require the following terminology.
A connected graph is \emph{biconnected} if it has at least 2 vertices and the removal of any vertex will always produce a connected graph.
We remark that the set of biconnected graphs is exactly the set of 2-connected graphs with the addition of the complete graph with two vertices.
Given a graph $G=(V,E)$, we say that a subset $F \subset E$ is a \emph{biconnected component} of $G$ if the subgraph $H=(V[F],F)$ induced by $F$ is biconnected and no subgraph of $G$ containing $H$ that is not $H$ itself is biconnected.
The biconnected components of a graph form a partition of its edge set.

\begin{proposition}\label{p:1d}
	Let $G$ be a graph and fix $k$ to be the number of biconnected components of $G$.
	Then $c_1(G) < \infty$ if and only if $G$ is connected.
	If $G$ is connected with at least one edge then $c_1(G) = 2^{k-1}$.
\end{proposition}

\begin{proof}
	If we glue two connected graphs $H_1$ and $H_2$ at exactly one vertex to form a graph $H$ then $c_1(H) = 2c_1(H_1)c_1(H_2)$.
	The result now follows from the observation that $c_1(G) = 1$ if $G$ is biconnected.
\end{proof}

As the number of biconnected components of a graph can be computed in linear time~\cite{HopcroftTarjan}, it follows that there exists a linear time algorithm for computing $c_1(G)$.
The current fastest deterministic algorithm for computing $c_2(G)$ when $G$ is minimally 2-rigid runs in exponential time~\cite{PlaneCount}.
There are no other known deterministic algorithms for computing $c_d(G)$ outside of restricting to specific families of graphs (\eg chordal graphs).




\section{Counting realisations on a complex sphere}\label{sec:countsphere}

Our aim in this section is to prove that the definition of the spherical $d$-realisation number alluded to in the introduction can be stated in a rigorous manner (see \Cref{def:countsphere}).

\subsection{Complex spherical rigidity map}

We define
\begin{align*}
	\mathbb{S}_{\mathbb{C}}^d := \left\{ x \in \mathbb{C}^{d+1} : \|x\|^2=1 \right\}
\end{align*}
to be the \emph{complexification of the $d$-dimensional (unit) sphere}.
As $\mathbb{S}_{\mathbb{C}}^d$ is a smooth connected variety (and hence irreducible),
so too is $(\mathbb{S}_{\mathbb{C}}^{d})^V$ for any finite set $V$.
Like how we chose to consider a larger set of realisations by switching from $(\mathbb{R}^d)^V$ to~$(\mathbb{C}^d)^V$,
we now consider spherical realisations in $\mathbb{S}_\mathbb{C}^d$ by extending the set $(\mathbb{S}^d)^V$ to the set $(\mathbb{S}_{\mathbb{C}}^d)^V$.
Any spherical realisation in $p \in (\mathbb{S}^d)^V$ is now said to be a \emph{real spherical realisation}.
At a spherical realisation $p \in (\mathbb{S}_{\mathbb{C}}^{d})^{V}$,
the linear space
\begin{align*}
	T_p (\mathbb{S}_{\mathbb{C}}^d)^V : =\left\{x \in (\mathbb{C}^{d+1})^V : x_v \cdot p_v=0 \text{ for all } v \in V \right\}
\end{align*}
is the tangent space of $(\mathbb{S}_{\mathbb{C}}^{d})^{V}$ at $p$.

For any graph $G=(V,E)$ we define the \emph{complex spherical rigidity map} to be the map
\begin{align*}
	s_{G,d} \colon (\mathbb{S}_{\mathbb{C}}^{d})^{V} \rightarrow \mathbb{C}^{E}, ~ p \mapsto \left( \frac{1}{2}\|p_v-p_w\|^2 \right)_{vw \in E} =\left( 1 - p_v \cdot p_w\right)_{vw \in E}.
\end{align*}
We denote the Zariski closure of the image of $s_{G,d}$ by $\ell^*_d(G)$.
Since the domain of $s_{G,d}$ is irreducible,
$\ell_d^*(G)$ is a variety.
The derivative of $s_{G,d}$ at a point $p \in (\mathbb{S}_{\mathbb{C}}^{d})^{V}$ can be seen to be the linear map
\begin{align*}
	\diff s_{G,d}(p) \colon T_p (\mathbb{S}_{\mathbb{C}}^{d})^{V} &\rightarrow \mathbb{C}^{E},\\
	x &\mapsto \left( (p_v-p_w) \cdot (x_v-x_w) \right)_{vw \in E} = \left(-(p_v \cdot x_w + p_w \cdot x_v) \right)_{vw \in E}.
\end{align*}
We define two realisations $p,q \in (\mathbb{S}_{\mathbb{C}}^{d})^V$ to be \emph{congruent} (denoted by $p \sim q$) if and only if $p = Aq$ for some $A \in O(d+1, \mathbb{C})$.
If the vertex set of $p$ linearly spans $\mathbb{C}^{d+1}$,
then two realisations $p,q$ are congruent if and only if $s_{K_V,d}(p) = s_{K_V,d}(q)$.
For all $p \in (S_\mathbb{C}^{d})^V$,
we define $C^*_d(G,p) := s^{-1}_{G,d} (s_{G,d}(p))/_\sim$ to be the \emph{spherical realisation space of $(G,p)$}.


We can link the infinitesimal properties of rigidity in $\mathbb{C}^d$ and in $\mathbb{S}^d_\mathbb{C}$ with the following result.

\begin{lemma}\label{l:equivrank}
	Let $p \in (\mathbb{C}^{d})^V$ and $q \in (\mathbb{S}_{\mathbb{C}}^{d})^V$ be such that $[q_{v}]_{d+1} \neq 0$ and $[p_v]_i = [q_v]_i/[q_v]_{d+1}$ for all $v \in V$ and $i \in \{1,\ldots,d\}$.
	Then $\rank \diff f_{G,d}(p)= \rank \diff s_{G,d}(q)$.
\end{lemma}

\begin{proof}
	Define the bijective linear map
	\begin{align*}
		\varphi \colon T_q (\mathbb{S}_{\mathbb{C}}^{d})^{V} \rightarrow (\mathbb{C}^{d})^V, ~ x \mapsto \left( \frac{[x_v]_1}{[q_v]_{d+1}},\ldots, \frac{[x_v]_d}{[q_v]_{d+1}} \right)_{v \in V}.
	\end{align*}
	For any $u \in T_q (\mathbb{S}_{\mathbb{C}}^{d})^{V}$ we have
	\begin{eqnarray*}
		&~& (p_v-p_w)\cdot (\varphi(u)_v- \varphi(u)_w)\\
		&=& \sum_{i=1}^d \frac{[q_v]_i[u_v]_i}{[q_v]^2_{d+1}} +\frac{[q_w]_i[u_w]_i}{[q_w]^2_{d+1}} - \frac{[q_v]_i[u_w]_i + [q_w]_i[u_v]_i}{[q_v]_{d+1}[q_w]_{d+1}} \\
		&=& -\frac{[q_v]_{d+1}[u_v]_{d+1}}{[q_v]^2_{d+1}}-\frac{[q_w]_{d+1}[u_w]_{d+1}}{[q_w]^2_{d+1}} \\
		&&+ \frac{[q_v]_{d+1}[u_w]_{d+1} +[q_w]_{d+1}[u_v]_{d+1}}{[q_v]_{d+1}[q_w]_{d+1}} - \frac{q_v\cdot u_w + q_w \cdot u_v}{[q_v]_{d+1}[q_w]_{d+1}} \\
		&=& - \frac{q_v \cdot u_w + q_w \cdot u_v}{[q_v]_{d+1}[q_w]_{d+1}}.
	\end{eqnarray*}
		It follows that $u \in \ker \diff s_{G,d}(q)$ if and only if $\varphi(u) \in \ker \diff f_{G,d}(p)$.
		As $\varphi$ is bijective,
		we have $\rank \diff f_{G,d}(p)= \rank \diff s_{G,d}(q)$ as required.
\end{proof}

From \Cref{l:equivrank} we know that a graph is rigid (respectively, independent) in $\mathbb{R}^d$ if and only if it is rigid (respectively, independent) in $\mathbb{S}^d$.
Due to this, we can obtain an analogue of \Cref{l:domind} for the map $s_{G,d}$.

\begin{lemma}\label{l:domindsphere}
	Let $G=(V,E)$ be any graph.
	Then the following are equivalent:
	\begin{enumerate}
		\item \label{l:sdomind1} $G$ is $d$-independent.
		\item \label{l:sdomind2} The map $s_{G,d}$ is dominant.	
		\item \label{l:sdomind3} $\ell_d^*(G) = \mathbb{C}^E$.
	\end{enumerate}
\end{lemma} 

\begin{proof}
	By the definition of a dominant map, \ref{l:sdomind2} and \ref{l:sdomind3} are equivalent.
	Hence, we only need to prove that \ref{l:sdomind1} and \ref{l:sdomind2} are equivalent.
	By \Cref{l:domind},
	we only need to show that $s_{G,d}$ is dominant if and only if $f_{G,d}$ is dominant.
	This is equivalent to proving the following:
	there exists $p \in (\mathbb{C}^{d})^V$ with $\rank \diff f_{G,d}(p) = |E|$ if and only if
	there exists a $q \in (\mathbb{S}_{\mathbb{C}}^{d})^V$ with $\rank \diff s_{G,d}(p) = |E|$ (see \Cref{borel91}).
	Hence, the result holds by \Cref{l:equivrank}.
\end{proof}


\subsection{Defining the spherical \texorpdfstring{$d$-realisation}{d-realisation} number}

Let $G=(V,E)$ be a graph with at least $d+1$ vertices and fix a sequence of $d$ vertices $v_1,\ldots, v_d$.
We now define the algebraic set
\begin{align}\label{eq:yset}
	Y_{G,d} := \left\{ p \in (\mathbb{S}^{d}_\mathbb{C})^V : [p_{v_k}]_j=0 \text{ if } j \geq k + 1 \text{ and }  [p_{v_1}]_1=1 \right\}
\end{align}
Since $Y_{G,d}$ is a connected smooth manifold,
it is irreducible.
We further note that $Y_{G,d}$ has dimension $d|V| - \binom{d+1}{2}$.
With this, we define the morphism
\begin{align*}
	\tilde{s}_{G,d} : Y_{G,d} \rightarrow \ell^*_d(G), ~ p \mapsto s_{G,d}(p),
\end{align*}
i.\,e., the restriction of $s_{G,d}$ to the domain $Y_{G,d}$ and the codomain $\ell^*_d(G)$.
Since every $d$-dimensional realisation of $G$ is equivalent to at least one element of $Y_{G,d}$ it follows that $\tilde{s}_{G,d}$ is a dominant map.

The methods utilised in \Cref{l:xgd} can easily be adapted to work for realisations in the $d$-dimensional sphere.

\begin{lemma}\label{l:ygd}
	Let $G=(V,E)$ be a graph with $|V| \geq d+1$,
	and fix a sequence of $d$ vertices $v_1,\ldots, v_d$.
	Then the image of $\tilde{s}_{G,d}$ is Zariski dense in the image of $s_{G,d}$ (and so $\tilde{s}_{G,d}$ is dominant),
	and
	\begin{align*}
		\left|\tilde{s}_{G,d}^{-1}\left(s_{G,d}(p) \right) \right| = 2^{d} |C^*_d(G,p)|
	\end{align*}
	for almost all $p \in (\mathbb{S}_\mathbb{C}^d)^V$.
\end{lemma}


The spherical version of \Cref{p:dom} is also proved in an analogous way.


\begin{proposition}\label{p:domsphere}
	Let $G=(V,E)$ be a graph with $|V| \geq d+1$.
	Then the following are equivalent:
	\begin{enumerate}
		\item\label{p:domsphere1} $G$ is $d$-rigid.
		\item\label{p:domsphere2} The map $\tilde{s}_{G,d}$ is generically finite.
		\item\label{p:domsphere3} There exists an $n \in \mathbb{N}$ and a non-empty Zariski open subset $U \subset (\mathbb{S}_\mathbb{C}^d)^V$ where $ |C^*_d(G,p)|=n$ for all $p \in U$.
		Furthermore,
		if $p \in U$ and $q$ is an equivalent $d$-dimensional spherical realisation of $G$,
		then $\rank \diff s_{G,d}(q) =  d|V| - \binom{d+1}{2}$.
	\end{enumerate}
\end{proposition} 

Using \Cref{p:domsphere} we can now make the following well-defined definition of the spherical $d$-realisation number for a minimally $d$-rigid graph.

\begin{definition}\label{def:countsphere}
	The \emph{spherical $d$-realisation number} of a graph $G=(V,E)$ is an element of $\mathbb{N} \cup \{\infty\}$ given by
	\begin{align*}
		c^*_d(G) :=		
		\begin{cases}
			|C^*_d(G,p)| \text{ for a general $p \in (\mathbb{S}_\mathbb{C}^d)^V$} &\text{if } |V| \geq d+1, \\
			 1 &\text{if } |V| \leq d \text{ and $G$ is complete}, \\
			 \infty &\text{if } |V| \leq d \text{ and $G$ is not complete}.
		\end{cases}
	\end{align*}
\end{definition}

It follows from \Cref{p:domsphere} that a graph $G$ is $d$-rigid if and only if $c_d(G) < \infty$.
It is worth noting that,
while $c^*_d(G) < \infty$ if and only if $c_d(G) < \infty$,
this does not mean the two numbers are always equal.
This can be seen in the following example.

\begin{example}
  We consider the 3-prism graph.
  It is minimally 2-rigid with six vertices, all of which have degree three.
  For this graph we have $c_2(G)=12$, which can be computed with the algorithm described in~\cite{PlaneCount}.
  In fact, there exist edge-length assignments for which we also obtain exactly 12 real realisations (see original Figure~2).
  However, the same graph embedded on the sphere gives $c_2^*(G)=16$,
  which can be computed by the algorithm described in~\cite{SphereCount}.
  Again there exist edge-length assignments such that each realisation is real (see original Figure~3).
  (Both algorithms use the alternative definitions for realisation numbers and so the outputted numbers first need to be halved.)
\end{example}

The natural analogue of \Cref{cor:realstuff} also holds for spheres by a similar method.

\begin{corollary}\label{cor:sphrealstuff}
	For each graph $G=(V,E)$,
	there exists a non-empty Zariski open subset $U_\mathbb{R} \subset (\mathbb{S}^d)^V$ of real $d$-dimensional spherical realisations $p$ where $|C^*_d(G,p)| = c^*_d(G)$.
\end{corollary}

\begin{remark}
	It is currently open whether every graph $G$ has a real generic $d$-dimensional spherical realisation that is equivalent (modulo congruences) to exactly $c_d^*(G)$ real spherical realisations.
	The counter-example used in the previous section for real realisations in the plane required the existence of a graph $H$ where $c_2(H)$ is odd.
	Similarly, any graph $G$ where $c^*_2(G)$ is odd would be a counter-example to the result, however no such graph has yet been found.
\end{remark}

We finish the section by observing that the spherical realisation number is always equal to the realisation number when $d=1$,
which can be proven using the same method as \Cref{p:1d}.

\begin{proposition}\label{p:1dsame}
	For any graph $G$ we have $c_1^*(G) =c_1(G)$.
\end{proposition}

\begin{remark}\label{rem:hyperbolic}
	All of the results in this section can be extended from the sphere to the light cone of any pseudo-Euclidean space.
	To be specific,
	we define for each $1 \leq k \leq d$ the quadratic form
	\begin{align*}
		q_{d,k}: \mathbb{C}^{d+1} \rightarrow \mathbb{C}, ~ (x_0,\ldots,x_{d}) \mapsto \sum_{j=0}^k x_j^2 - \sum_{j=k+1}^d x_j^2,
	\end{align*}
	and we define the light cone
	\begin{align*}
		\mathbb{H}_{d,k} := \left\{ x \in \mathbb{C}^{d+1} : q_{d,k}(x) = 0 \right\}.
	\end{align*}
	The map
	\begin{align*}
		T:\mathbb{H}_{d,k} \rightarrow \mathbb{H}_{d,d}, ~ (x_0,\ldots,x_d) \mapsto (x_0,\ldots, x_k, ix_{k+1},\ldots,x_d)
	\end{align*}
	is an isometry (in the sense that $q_{d,d}(T(x) - T(y)) = q_{d,k}(x-y)$),
	and the space $\mathbb{H}_{d,d}$ is isomorphic to $\mathbb{S}_{\mathbb{C}}^d$.
	Hence, the realisation number of a graph $G$ in any space $\mathbb{H}_{d,k}$ is equal to $c^*_d(G)$.
	Notably,
	the hyperbolic realisation number of a graph (equivalently, the space $\mathbb{H}_{d,d-1}$) is equal to the spherical realisation number of a graph. 
\end{remark}



\section{Coning and the spherical realisation number}\label{sec:cone}

In this section we prove \Cref{t:conecount} by observing that the number of realisations of a framework is projectively invariant (\Cref{kl:proj}).
We then prove that $c_d(G*o) = c_d^*(G*o)$ for any coned graph $G*o$ (original Theorem~5.2).

\subsection{Proof of \texorpdfstring{\Cref{t:conecount}}{Theorem 1.2}}

Throughout this subsection we fix the vertices used to define our various varieties in the following way.
Let $G=(V,E)$ be a graph with $|V| \geq d+1$ and let $G*o$ be its cone.
We fix the vertices $v_1,\ldots, v_d \in V$ to define the set $X_{G,d}$ as given in \eqref{eq:xset}.
From this, we fix the set
\begin{align*}
	X_{G*o,d+1} = \left\{ p' \in (\mathbb{C}^{d+1})^{V*o} : p'_o = \mathbf{0}, ~  [p'_{v_k}]_{j+1}=0 \text{ if } k \leq j \leq d  \right\}\,.
\end{align*}
(This is equivalent to defining the set $X_{G*o,d+1}$ as in \eqref{eq:xset} with the vertices $w_1,\ldots,  w_{d+1}$ by setting $w_1 = o$ and $w_{j+1} = v_j$ for all $j =1,\ldots, d$.)
We also assume that the fixed vertices $v_1,\ldots,v_d$ used to define $X_{G,d}$ and $Y_{G,d}$ (see \eqref{eq:yset}) are always identical.



\begin{lemma}\label{kl:proj}
	Given a graph $G=(V,E)$ with $|V| \geq d+1$,
	choose any $p \in Y_{G,d}$ and $(r_v)_{v \in V} \in (\mathbb{C}\setminus \{0\})^{V}$.
	Define $p' \in X_{G*o,d+1}$ with $p'_v = (r_vp_v,0)$ for each $v \in V$ and $p'_o=\mathbf{0}$.
	Then
	\begin{align*}
		\left| \tilde{f}_{G*o,d+1}^{-1}\left(\tilde{f}_{G*o,d+1}(p') \right) \right| = 2 \left| \tilde{s}_{G,d}^{-1}\Big(\tilde{s}_{G,d}(p) \Big) \right|.
	\end{align*} 
\end{lemma}

\begin{proof}
	Define $\lambda := \tilde{s}_{G,d}(p)=\left( 1 - p_v \cdot p_w\right)_{vw\in E}$ and $\lambda' := \tilde{f}_{G*o,d+1}(p')$.
	We first note that for any $vw \in E$ we have
	\begin{align}\label{eq:lambda}
		\lambda_{vw}' &= \frac{1}{2}\|p'_v- p'_w\|^2 \\
		\notag &= \frac{1}{2} \|p'_v\|^2 + \frac{1}{2} \|p'_w\|^2 - p'_v \cdot p'_w = \lambda'_{ov} + \lambda'_{ow} + (r_vr_w)(\lambda_{vw}-1).
	\end{align}
	Define $S \subset \tilde{f}_{G*o,d+1}^{-1}(\lambda')$ to be the set of realisations $q' \in X_{G*o,d+1}$ where $[q'_{v_1}]_1=r_{v_1}$.
	Note that for each $q' \in \tilde{f}_{G*o,d+1}^{-1}(\lambda')$ we have that
	\begin{align*}
		[q'_{v_1}]_1^2 = \|q'_{v_1} \|^2 = \|q'_{v_1} -q'_o \|^2 = \|p'_{v_1} -p'_o \|^2 = \|r_{v_1} p_{v_1}\|^2 = r_{v_1}^2;
	\end{align*}
	the first equality follows from $q' \in X_{G*o,d+1}$,
	the second from $ov_1 \in E*o$,
	the third from $ \tilde{f}_{G*o,d+1}(q' ) = \tilde{f}_{G*o,d+1}(p')$,
	the fourth from the construction of $p'$ from $p$,
	and the last from $p \in Y_{G,d}$.
	Hence, if $q' \in \tilde{f}_{G*o,d+1}^{-1}(\lambda') \setminus S$ then we have $[q'_{v_1}]_1 = -r_{v_1}$,
	and so $-q' \in S$,
	since 
	As $S \cap (-S) = \emptyset$ and $S \cup (-S) = \tilde{f}_{G*o,d+1}^{-1}(\lambda')$,
	we have that $|\tilde{f}_{G*o,d+1}^{-1}(\lambda')| = 2|S|$ (note that this cardinality need not be finite).
	It now suffices to prove that $|\tilde{s}_{G,d}^{-1}(\lambda)| = |S|$.
	
	Choose any $q \in \tilde{s}_{G,d}^{-1}(\lambda)$ and define $q'\in X_{G*o,d+1}$ by setting $q'_o=\mathbf{0}$ and $q'_v = r_v q_v$ for each $v \in V$.
	It is immediate that $[q'_{v_1}]_1 = r_{v_1}$ and $\frac{1}{2}\|q'_v\|^2 = \lambda'_{ov}$.
	For each $vw \in E$ we have
	\begin{align*}
		 \frac{1}{2}\|q'_v- q'_w\|^2 &=  \frac{1}{2}\|q'_v\|^2 +  \frac{1}{2}\|q'_w\|^2 - q'_v\cdot q'_w \\
		&= \lambda'_{ov} + \lambda'_{ow} + (r_vr_w)(\lambda_{vw}-1) \\
		&\stackrel{_\eqref{eq:lambda}}{=} \lambda'_{vw},
	\end{align*}
	and so $q' \in S$.
	Hence, $|\tilde{s}_{G,d}^{-1}(\lambda)| \leq |S|$.
	
	Now choose any $q' \in S$ and define $q \in (\mathbb{C}^{d+1})^V$ by setting $q_v = q'_v/r_v$ for each $v\in V$.
	We first note that $\|q_v\|^2 = \|q'_v - q'_o\|^2/r^2_v = 1$ for each $v \in V$,
	hence $q \in Y_{G,d}$.
	For each $vw \in E$ we have
	\begin{align*}
		q_v\cdot q_w &= \frac{q'_v\cdot q'_w}{r_vr_w} \\
		&= \frac{\|q'_v\|^2 + \|q'_w\|^2 - \|q'_v-q'_w\|^2}{2 r_vr_w} \\
		&= \frac{\lambda'_{ov} + \lambda'_{ow} - \lambda'_{vw}}{r_vr_w} \\
		&\stackrel{_\eqref{eq:lambda}}{=} 1 - \lambda_{vw},
	\end{align*}
	and so $q \in \tilde{s}_{G,d}^{-1}(\lambda)$.
	Hence, $|S |\leq |\tilde{s}_{G,d}^{-1}(\lambda)|$.
\end{proof}

A nice corollary of \Cref{kl:proj} is that the set $C_{d+1}(G*o,p)$ is projectively invariant,
in the sense that projecting each point $p_v$ along the complex line through the points $\{p_v,p_o\}$ by a non-zero complex scalar does not alter the number of equivalent realisations.

We are now ready to prove our first main theorem.

\begin{proof}[Proof of \Cref{t:conecount}]
	If $G=(V,E)$ is not $d$-rigid then the result follows from \Cref{t:cone}.
	If $|V| \leq d$ and $G$ is complete then $|V*o| \leq d+1$ and $G*o$ is complete,
	and so $c_{d+1}(G*o) = c^*_d(G)$.
	Suppose that $G=(V,E)$ is $d$-rigid with $|V| \geq d+1$.	
	As $G$ is $d$-rigid, the coned graph $G*o$ is $(d+1)$-rigid by \Cref{t:cone}.
	By \Cref{l:xgd} and \Cref{p:dom},
	there exists a Zariski open set $U' \subset (\mathbb{C}^{d+1})^{V*o}$ such that
	\begin{align}\label{e:conenumber}
		\left| \tilde{f}_{G*o,d+1}^{-1} \left(\tilde{f}_{G*o,d+1}(p') \right) \right| = 2^{d+1} c_{d+1}(G*o)
	\end{align}
	for all $p' \in U'$.
	Similarly, it follows from \Cref{l:ygd} and \Cref{p:domsphere} that there exists a Zariski open set $U \subset (\mathbb{S}_\mathbb{C})^{d}$ such that
	\begin{align}\label{e:spherenumber}
		\left| \tilde{s}_{G,d}^{-1} \Big(\tilde{s}_{G,d}(p) \Big) \right| = 2^{d} c_{d}^*(G)
	\end{align}
	for all $p \in U$.
	Fix $U^* := U \times (\mathbb{C} \setminus \{0\})^V$.
	Since $U^*$ is a Zariski open subset of the variety $Y_{G,d} \times \mathbb{C}^V$,
	it is an open and dense subset of $Y_{G,d} \times \mathbb{C}^V$ with respect to the metric topology.
	Similarly,
	since $U'$ is a Zariski open subset of the variety $ X_{G *o,d+1}$,
	it is an open and dense subset of $X_{G *o,d+1}$ with respect to the metric topology.
	
	Define the surjective (and hence dominant) morphism 
	\begin{align*}
		\phi : Y_{G,d} \times \mathbb{C}^V \rightarrow X_{G *o,d+1},
	\end{align*}
	where, given $p'= \phi(p,r)$,
	we have $p'_v=r_vp_v$ for all $v \in V$ and $p'_o=\mathbf{0}$.
	As $\phi$ is dominant and $\dim Y_{G,d} \times \mathbb{C}^V = \dim X_{G *o,d+1}$,
	it follows from \Cref{t:deg} and the inverse mapping theorem for holomorphic maps (see~\cite[Theorem 7.5]{FG02}) that the image of any open dense subset of $Y_{G,d} \times \mathbb{C}^V$ (with respect to the metric topology) contains an open subset of $X_{G*o,d+1}$ (with respect to the metric topology).
	Hence, the set $\phi( U^*)$ has a non-empty interior with respect to the metric topology of $X_{G*o,d+1}$.
	Since $U'$ is an open dense subset of $X_{G*o,d+1}$ with respect to the metric topology,
	$\phi( U^*) \cap U'$ contains a realisation $p'$.
	Choose $p \in U$ and $r \in (\mathbb{C}\setminus \{0\})^V$ such that $\phi(p,r) = p'$.
	By \eqref{e:conenumber}, \eqref{e:spherenumber} and \Cref{kl:proj},
	we have that
	\begin{align*}
		2^{d+1} c_{d+1}(G*o) =\left| \tilde{f}_{G*o,d+1}^{-1}\left(\tilde{f}_{G*o,d+1}(p') \right) \right| = 2 \left| \tilde{s}_{G,d}^{-1}\Big(\tilde{s}_{G,d}(p) \Big) \right|= 2^{d+1} c^*_{d}(G),
	\end{align*}
	and so $c^*_{d}(G) = c_{d+1}(G*o)$.
\end{proof}




\appendix

\section{Dominant morphisms}
\label{sec:dom}

Dominant morphisms can be defined in a variety of different but equivalent ways.

\begin{theorem}[{\cite[Section AG, Theorem 17.3]{borel91}}]\label{borel91}
	Let $X \subset \mathbb{C}^m$ be an algebraic set and $Y \subset \mathbb{C}^n$ be a variety.
	Then the following are equivalent for any morphism $f :X \rightarrow Y$:
	\begin{enumerate}
		\item $f$ is dominant.
		\item For some irreducible component $X'$ of $X$, there exists a point $x \in X'$ such that $x$ is a non-singular point of $X'$ and $\rank \diff f(x) = \dim Y$.
		\item There exists a Zariski open subset $U \subset X$ where for each $x \in U$, $x$ is a non-singular point of $X$ and $\rank \diff f(x) = \dim Y$.
	\end{enumerate}
\end{theorem}

It follows immediately from \Cref{borel91} that for any varieties $X,Y$,
the existence of a dominant morphism from $X$ to $Y$ implies $\dim X \geq \dim Y$.
As can be seen by \Cref{t:deg},
more powerful statements can be attained relating to $f$ if $\dim X=\dim Y$.
To prove \Cref{t:deg}, we require the following two results.

\begin{corollary}[{\cite[Section 8, Corollary 1]{mumford}}]\label{mumford}
	Let $X \subset \mathbb{C}^m$ and $Y \subset \mathbb{C}^n$ be varieties and $f:X \rightarrow Y$ be a dominant morphism.
	Then there exists a non-empty Zariski open subset $U \subset Y$ such that $U \subset f(X)$, and for every $y \in U$, every irreducible component of the algebraic set $f^{-1}(y)$ has dimension $\dim X - \dim Y$.
\end{corollary}

\begin{corollary}[{\cite[Proposition 7.16]{Harris}}]\label{c:harris}
	Let $X \subset \mathbb{C}^m$ and $Y \subset \mathbb{C}^n$ be varieties and $f:X \rightarrow Y$ be a dominant morphism.
	Suppose that there exists a non-empty Zariski open subset $U \subset Y$ where $|f^{-1}(y)|<\infty$ for every $y \in Y$.
	Then there exists a $k \in \mathbb{N}$ and a non-empty Zariski open subset $U' \subset U$ where $|f^{-1}(y)| = k$ for every $y \in U'$.
\end{corollary}

\begin{proof}[Proof of \Cref{t:deg}]
	It is immediate that \ref{t:deg3} implies \ref{t:deg2}.
	Fix $U \subset Y$ to be the Zariski open set from \Cref{mumford}.
	An algebraic set is a finite set if and only if it is zero dimensional.
	Hence, \ref{t:deg1} and \ref{t:deg2} are equivalent.
	Suppose that $\dim X=\dim Y$.
	Fix $U' \subset U$ to be the non-empty Zariski open subset from \Cref{c:harris}.
	Define the set
	\begin{align*}
		C := \{x \in X : x \text{ is singular or } \rank \diff f(x) <n\}.
	\end{align*}
	By \Cref{borel91},
	$C$ is a proper Zariski closed subset of $X$.
	As $\dim C < \dim X = \dim Y$,
	the set $f(C)$ is not dense in $Y$.
	It now follows that the complement of the Zariski closure of $f(C)$ in $U$ is a non-empty Zariski open subset of $Y$.
	Hence, \ref{t:deg1} implies \ref{t:deg3},
	concluding the proof.
\end{proof}







\section{Proof of \texorpdfstring{\Cref{l:xgd}}{Lemma}}
\label{sec:xgd}

We begin with the following technical lemma describing when we can ``rotate'' a framework into our set $X_{G,d}$.

\begin{lemma}\label{l:rotate}
	Let $G=(V,E)$ be a graph with $|V| \geq d+1$,
	and fix a sequence of $d$ vertices $v_1,\ldots, v_d$.
	For each realisation $p \in (\mathbb{C}^d)^V$,
	define the $(d-1) \times (d-1)$ symmetric matrix
	\begin{align*}
		\mathbb{G}(p) := 		
		\begin{bmatrix}
			(p_{v_2} - p_{v_1})^T \\
			\vdots \\
			(p_{v_d} - p_{v_1})^T
		\end{bmatrix}
		\begin{bmatrix}
			p_{v_2} - p_{v_1} & \cdots & p_{v_d} - p_{v_1}
		\end{bmatrix}.
	\end{align*}
	Then there exists a realisation $q \in X_{G,d}$ congruent to $p$ if $\mathbb{G}(p)$ only has non-zero leading principal minors.\footnote{A \emph{leading principal minor (of order $n$)} of a square matrix is the determinant of the matrix formed by taking the first $n$ rows and columns.}
\end{lemma}

The conditions stated in \Cref{l:rotate} seem rather bizarre,
especially since no such conditions are required if we restrict ourselves to real realisations.
It is also rather easy to see they are not necessary either:
simply choose any realisation $p \in X_{G,d}$ where $p_{v_i} = \mathbf{0}$ for each $i \in \{1,\ldots,d\}$.
To see why we need to be so cautious,
take $G=(V,E)$ to be the complete graph with 4 vertices $v_1,v_2,v_3,v_4$,
and let $p$ be the 3-dimensional realisation of $G$ where
\begin{align*}
	p_{v_1} = (0,0,0), \qquad p_{v_2} = (1,0,0), \qquad p_{v_3} = (2,1,i), \qquad p_{v_4} = (0,1,0).
\end{align*}
In this specific case, there are no realisations in $X_{G,d}$ that are congruent to $p$.
The framework $(G,p)$ also has the interesting properties that its vertices affinely span $\mathbb{C}^3$ and $\|p_v-p_w\|^2 >0$ for every edge $vw \in E$.
Since the matrix
\begin{align*}
	\mathbb{G}(p) =
	\begin{bmatrix}
		1 & 0 & 0 \\
		2 & 1 & i
	\end{bmatrix}
	\begin{bmatrix}
		1 & 2 \\
		0 & 1 \\
		0 & i
	\end{bmatrix}
	=
	\begin{bmatrix}
		1 & 2 \\
		2 & 4
	\end{bmatrix}
\end{align*}
has rank 1,
our constructed realisation is not a counter-example to \Cref{l:rotate}.

\begin{proof}
	Fix $p \in (\mathbb{C}^d)^V$ with the stipulated properties.
	With this,
	define $p^1 \in (\mathbb{C}^d)^V$ to be the realisation where $p^1_v = p_v - p_{v_1}$ for each $v \in V$.
	If $d=1$ then $p^1 \in X_{G,1}$,
	so suppose that $d \geq 2$.
	We now observe that the matrix $\mathbb{G}(p^1) = \mathbb{G}(p)$.
	In fact, a much stronger statement is true:
	if $q$ is a congruent realisation then $\mathbb{G}(q) = \mathbb{G}(p)$.
	
	We now form the sequence of congruent realisations $p^1,\ldots,p^d$ in the following inductive way.
	Fix $n \in \{2,\ldots,d\}$,
	and suppose that $p^1,\ldots,p^{n-1}$ have already been constructed such that $[p^{n-1}_{v_k}]_j =0$ if $k \leq \min \{j, n-1\}$.
	For each $k \in \{1,\ldots,d\}$,
	fix $e_k$ to be the vector with $[e_k]_j = 1$ if $j=k$ and $[e_k]_j = 0$ otherwise.
	We observe here that the linear space spanned by the vectors $e_1,\ldots, e_{n-2}$ contains the points $p^{n-1}_{v_2},\ldots, p^{n-1}_{v_{n-1}}$.
	Fix $z$ to be the vector formed from $p^{n-1}_{v_{n}}$ by replacing its first $n-2$ coordinates with zeroes.
	It is immediate that $z \cdot e_j = 0$ for all $j \in \{1 , \ldots, n-2\}$.
	
	Suppose for contradiction that $\|z\|^2 = 0$,
	and fix $y = p^{n-1}_{v_{n}} - z$.
	Set $A = [p^{n-1}_{v_2} ~ \cdots ~ p^{n-1}_{v_{n-1}}]$.
	The determinant of $A^T A$ is the leading principal minor of order $n-2$ of $\mathbb{G}(p) = \mathbb{G}(p^{n-1})$,
	and so is non-zero.
	Hence, the column span of $A$ is exactly the linear space spanned by $e_1,\ldots, e_{n-2}$;
	importantly,
	this implies that $y$ is contained in the column span of $A$.
	The leading principal minor of order $n-1$ of $\mathbb{G}(p) = \mathbb{G}(p^{n-1})$ is the determinant of the $(n-1) \times (n-1)$ matrix
	\begin{align*}
		B &=		
		\begin{bmatrix}
			(p^{n-1}_{v_2})^T \\
			\vdots \\
			(p^{n-1}_{v_{n}})^T
		\end{bmatrix}
		\begin{bmatrix}
			p^{n-1}_{v_2} & \cdots & p^{n-1}_{v_{n}}
		\end{bmatrix}\\
		&=
		\begin{bmatrix}
			A^T A & A^T (y+z) \\
			(y+z)^T A & (y+z)^T (y+z)
		\end{bmatrix}\\
		&=
		\begin{bmatrix}
			A^T A & A^T y \\
			y^T A & y^T y
		\end{bmatrix}\\
		&= 
		\begin{bmatrix}
			A^T \\
			y^T
		\end{bmatrix}
		\begin{bmatrix}
			A & y
		\end{bmatrix}.
	\end{align*}
	As $y$ is contained in the column span of $A$,
	the matrix $[A ~ y]$ has rank at most $n-2$.
	However, this implies $\det B = 0$,
	contradicting our leading principal minor assumption for $\mathbb{G}(p)$.
	Hence, $\|z\|^2 \neq 0$.
	
	Fix $x := z/\|z\|^2$.
	Define the $d \times d$ matrix $\tilde{M}$ where $\tilde{M}^T = [ e_1 ~ \cdots ~ e_{n-2} ~ x ~ \mathbf{0} ~ \cdots ~ \mathbf{0}]$.
	We note that the matrix $\tilde{M}$ is an isometry (in the quadratic form sense) from the linear subspace of $\mathbb{C}^d$ spanned by $p^{n-1}_{v_2},\ldots, p^{n-1}_{v_{n}}$ to the linear subspace $\mathbb{C}^{n-1} \times \{0\}^{d-n+1}$.
	By Witt's theorem (see, for example,~\cite[Theorem 11.15]{roman}),
	there exists a matrix $M \in O(d, \mathbb{C})$ where $Mp^{n-1}_{v_j} = \tilde{M}p^{n-1}_{v_j}$ for each $j \in \{1,\ldots,n\}$.
	With this we now fix $p^n_v = Mp^{n-1}_v$ for each $v \in V$.
	The result now follows from completing the inductive argument and fixing $q = p^d$.
\end{proof}


With this, we can prove \Cref{l:xgd}.


\begin{proof}[Proof of \Cref{l:xgd}]
	For every $p \in (\mathbb{C}^d)^V$,
	let $\mathbb{G}(p)$ be the matrix defined in the statement of \Cref{l:rotate}.
	With this, we fix the proper algebraic set
	\begin{align*}
		Z := \left\{ p \in (\mathbb{C}^d)^V :
			\begin{array}{l}
				\text{$\mathbb{G}(p)$ has a zero leading principal minor,}\\
				\text{or $p$ does not affinely span $\mathbb{C}^d$}
			\end{array}
		 \right\}.
	\end{align*}
	Since $Z$ contains hypersurfaces in $(\mathbb{C}^d)^V$ and $Z \neq (\mathbb{C}^d)^V$
	it has dimension $d|V|-1$.
	By \Cref{l:rotate},
	the set $f_{G,d}((\mathbb{C}^d)^V \setminus Z)$ is a subset of the image of $\tilde{f}_{G,d}$.	
	Hence, the image of $\tilde{f}_{G,d}$ is Zariski dense in the image of $f_{G,d}$.
	
	Since $\dim \ell_d(G) = \rank \diff f_{G,d}(p)$ for any general realisation $p\in (\mathbb{C}^d)^V$,
	it follows from \Cref{t:asimowroth} that $G$ is $d$-rigid if and only if $\dim \ell_d(G) = d|V|- \binom{d+1}{2}$.
	Suppose that $G$ is not $d$-rigid.
	The extension of \Cref{t:asimowroth} to complex realisations gives that the set $C_d(G,p)$ contains infinitely many points for almost all realisations $p$;
	in particular,
	it can be used to prove that for almost all $p \in (\mathbb{C}^d)^V \setminus Z$,
	there exists a sequence $(p^n)_{n \in \mathbb{N}}$ of realisations where $p^n$ is not congruent to either $p^m$ or $p$ for each $m \neq n$, and $p^n \rightarrow p$ as $n \rightarrow \infty$ in the standard complex norm for $(\mathbb{C}^d)^V$ (note that this is an actual metric and not the quadratic form $\| \cdot\|^2$).
	An application of \Cref{l:rotate} to both $p$ and each $p^n$ gives that $\tilde{f}_{G,d}^{-1}\left(f_{G,d}(p)\right)$ also contains infinitely many points as is required.
	
	We now suppose that $G$ is $d$-rigid,
	i.\,e., $\dim \ell_d(G) = d|V|- \binom{d+1}{2}$.
	We split the proof into three cases.
	
	\textbf{(Case 1: $f_{G,d}(Z)$ is Zariski dense in $\ell_d(G)$ and $|V| \geq d+2$.)}
	In this case,
	the restriction of the map $f_{G,d}$ to $Z$ and $\ell_d(G)$, which we now denote by $g : Z \rightarrow \ell_d(G)$, is dominant.
	Hence, by \Cref{borel91},
	there exists a non-empty Zariski open subset $U \subset Z$ where 
	\begin{align*}
		\rank \diff g(p) = \dim \ell_d(G) = d|V|- \binom{d+1}{2}
	\end{align*}
	for each $p \in U$.
	Since $\rank \diff g(p) \leq \rank \diff f_{G,d}(p)$,
	we have that $\rank \diff f_{G,d}(p) = d|V| - \binom{d+1}{2}$ for each $p \in U$.
	As $Z$ has dimension $d|V|-1$,
	it contains an irreducible component of dimension $d|V|-1$.
	It follows that the nullity of $\diff g(p)$ is $\binom{d+1}{2} - 1$ for each $p \in U$.	
	If $q$ is congruent to $p \in Z$ then $\mathbb{G}(q) = \mathbb{G}(p)$ and $q$ is affinely spanning if and only if $p$ is,
	hence $Z$ is closed under congruences.
	This implies that for each $p \in U$,
	the map from the affine congruences of $\mathbb{C}^d$ to the corresponding congruent realisation of $p$ in $Z$ is not injective,
	as otherwise the nullity of $\diff g(p)$ would be at least $\binom{d+1}{2}$ for each $p \in U$.
	Hence, any realisation in $U$ cannot affinely span $\mathbb{C}^d$.
	The set of realisations of $G$ that do not affinely span $\mathbb{C}^d$ is the intersection of $|V| - d \geq 2$ pairwise-distinct hypersurfaces,
	and so $U$ is contained in a $(d|V|-2)$-dimensional subvariety in $Z$.
	However, this now contradicts that $U$ is a non-empty Zariski open subset of the $(d|V|-1)$-dimensional algebraic set $Z$.
	
	
	\textbf{(Case 2: $f_{G,d}(Z)$ is not Zariski dense in $\ell_d(G)$ and $|V| \geq d+2$.)}
	Fix the non-empty Zariski open set
	\begin{align*}
		X := \left\{ p \in X_{G,d} : \tilde{f}_{G,d}(p) \text{ is not contained in the closure of $f_{G,d}(Z)$ }   \right\}.
	\end{align*}
	Choose any $p \in X$.
	By \Cref{l:rotate}, for each realisation $p' \in (\mathbb{C}^d)^V$ that is equivalent to $p$,
	there exists some other realisation $p'' \in X_{G,d}$ that is congruent to $p'$.
	Hence, it now suffices to prove that $p$ is congruent to exactly $2^d$ realisations (including itself) in $X_{G,d}$.
	
	Choose any $q \in X_{G,d}$ congruent to $p$,
	and let $A \in O(d,\mathbb{C})$ be the matrix that maps the vertices of $p$ to the vertices of $q$.
	As $p,q \in X_{G,d}$, we have 
	\begin{align*}
		\sum_{n=1}^k A_{j,n} [p_{v_{k+1}}]_n = [q_{v_{k+1}}]_j = 0
	\end{align*}
	for each $1 \leq k < j \leq d$,
	hence $A_{j,k}=0$ for each $1 \leq k < j \leq d$.
	Since $A^T$ maps the vertices of $q$ to the vertices of $p$,
	we similarly have $A_{j,k}=0$ for each $1 \leq j < k \leq d$.
	As the vertices of $p$ linearly span $\mathbb{C}^d$,
	the set $\tilde{f}_{G,d}^{-1}\left(f_{G,d}(p) \right)$ is in one-to-one correspondence with the diagonal orthogonal matrices.
	The set of diagonal orthogonal matrices is exactly the set of diagonal matrices $A$ where $A_{j,j} = \pm 1$ for each $j \in \{1,\ldots, d\}$,
	hence there are $2^d$ of them.
	This now concludes the proof.
	
	\textbf{(Case 3: $|V| =d+1$.)}	
	Since $G$ is $d$-rigid, it can be easily verified that $G$ must be the complete graph on $d+1$ vertices.
	Furthermore,
	every element $p \in (\mathbb{C}^d)^V \setminus Z$ is only equivalent to congruent realisations (\cite[Corollary 8]{Gortler2014}),
	and each equivalent realisation is also contained in $(\mathbb{C}^d)^V \setminus Z$.
	We now define $X = X_{G,d} \setminus Z$ and repeat the same technique as was utilised in Case~2.
	This now concludes the proof.
\end{proof}





\begin{thebibliography}{MMMMMM}
\bibitem[1]{Abbot} Abbott, Timothy G. Generalizations of Kempe’s universality theorem, Masters thesis, Massachusetts Institute of Technology (2008) \href{https://alco.centre-mersenne.org/articles/10.5802/alco.390/#Abbot}{Publisher reference}.
\bibitem[2]{AsimowRothI} Asimow, L.; Roth, B. The rigidity of graphs, Trans. Amer. Math. Soc., Volume 245 (1978), pp. 279-289 \href{https://alco.centre-mersenne.org/articles/10.5802/alco.390/#AsimowRothI}{Publisher reference}.
\bibitem[3]{Bartzos2021} Bartzos, Evangelos; Emiris, Ioannis Z.; Legerský, Jan; Tsigaridas, Elias On the maximal number of real embeddings of minimally rigid graphs in ℝ 2 , ℝ 3  and S 2 , J. Symbolic Comput., Volume 102 (2021), pp. 189-208 \href{https://alco.centre-mersenne.org/articles/10.5802/alco.390/#Bartzos2021}{Publisher reference}.
\bibitem[4]{Bartzos2020} Bartzos, Evangelos; Emiris, Ioannis Z.; Schicho, Josef On the multihomogeneous Bézout bound on the number of embeddings of minimally rigid graphs, Appl. Algebra Engrg. Comm. Comput., Volume 31 (2020) no. 5-6, pp. 325-357 \href{https://alco.centre-mersenne.org/articles/10.5802/alco.390/#Bartzos2020}{Publisher reference}.
\bibitem[5]{Bartzos2023} Bartzos, Evangelos; Emiris, Ioannis Z.; Tzamos, Charalambos An asymptotic upper bound for graph embeddings, Discrete Appl. Math., Volume 327 (2023), pp. 157-177 \href{https://alco.centre-mersenne.org/articles/10.5802/alco.390/#Bartzos2023}{Publisher reference}.
\bibitem[6]{BartzosUpperBound} Bartzos, Evangelos; Emiris, Ioannis Z.; Vidunas, Raimundas New upper bounds for the number of embeddings of minimally rigid graphs, Discrete Comput. Geom., Volume 68 (2022) no. 3, pp. 796-816 \href{https://alco.centre-mersenne.org/articles/10.5802/alco.390/#BartzosUpperBound}{Publisher reference}.
\bibitem[7]{BorceaStreinu} Borcea, Ciprian; Streinu, Ileana The number of embeddings of minimally rigid graphs, Discrete Comput. Geom., Volume 31 (2004) no. 2, pp. 287-303 \href{https://alco.centre-mersenne.org/articles/10.5802/alco.390/#BorceaStreinu}{Publisher reference}.
\bibitem[8]{borel91} Borel, Armand Linear algebraic groups, Graduate Texts in Mathematics, 126, Springer-Verlag, New York, 1991, xii+288 pages \href{https://alco.centre-mersenne.org/articles/10.5802/alco.390/#borel91}{Publisher reference}.
\bibitem[9]{ZenodoAlg} Capco, Jose; Gallet, Matteo; Grasegger, Georg; Koutschan, Christoph; Lubbes, Niels; Schicho, Josef An algorithm for computing the number of realizations of a Laman graph, Zenodo, 2018 \href{https://alco.centre-mersenne.org/articles/10.5802/alco.390/#ZenodoAlg}{Publisher reference}.
\bibitem[10]{PlaneCount} Capco, Jose; Gallet, Matteo; Grasegger, Georg; Koutschan, Christoph; Lubbes, Niels; Schicho, Josef The number of realizations of a Laman graph, SIAM J. Appl. Algebra Geom., Volume 2 (2018) no. 1, pp. 94-125 \href{https://alco.centre-mersenne.org/articles/10.5802/alco.390/#PlaneCount}{Publisher reference}.
\bibitem[11]{ZenodoData} Capco, Jose; Gallet, Matteo; Grasegger, Georg; Koutschan, Christoph; Lubbes, Niels; Schicho, Josef The number of realizations of all Laman graphs with at most 12 vertices, Zenodo, 2018 \href{https://alco.centre-mersenne.org/articles/10.5802/alco.390/#ZenodoData}{Publisher reference}.
\bibitem[12]{cw10} Connelly, R.; Whiteley, W. J. Global rigidity: the effect of coning, Discrete Comput. Geom., Volume 43 (2010) no. 4, pp. 717-735 \href{https://alco.centre-mersenne.org/articles/10.5802/alco.390/#cw10}{Publisher reference}.
\bibitem[13]{EJNSTW19} Eftekhari, Yaser; Jackson, Bill; Nixon, Anthony; Schulze, Bernd; Tanigawa, Shin-ichi; Whiteley, Walter Point-hyperplane frameworks, slider joints, and rigidity preserving transformations, J. Combin. Theory Ser. B, Volume 135 (2019), pp. 44-74 \href{https://alco.centre-mersenne.org/articles/10.5802/alco.390/#EJNSTW19}{Publisher reference}.
\bibitem[14]{FG02} Fritzsche, Klaus; Grauert, Hans From holomorphic functions to complex manifolds, Graduate Texts in Mathematics, 213, Springer-Verlag, New York, 2002, xvi+392 pages \href{https://alco.centre-mersenne.org/articles/10.5802/alco.390/#FG02}{Publisher reference}.
\bibitem[15]{SphereCount} Gallet, Matteo; Grasegger, Georg; Schicho, Josef Counting realizations of Laman graphs on the sphere, Electron. J. Combin., Volume 27 (2020) no. 2, Paper no. 2.5, 18 pages \href{https://alco.centre-mersenne.org/articles/10.5802/alco.390/#SphereCount}{Publisher reference}.
\bibitem[16]{SphereAlg} Gallet, Matteo; Grasegger, Georg; Schicho, Josef Software for counting realizations of minimally rigid graphs on the sphere, Zenodo, 2022 \href{https://alco.centre-mersenne.org/articles/10.5802/alco.390/#SphereAlg}{Publisher reference}.
\bibitem[17]{Gortler2014} Gortler, Steven J.; Thurston, Dylan P. Generic global rigidity in complex and pseudo-Euclidean spaces, Rigidity and symmetry (Fields Inst. Commun.), Volume 70, Springer, New York, 2014, pp. 131-154 \href{https://alco.centre-mersenne.org/articles/10.5802/alco.390/#Gortler2014}{Publisher reference}.
\bibitem[18]{RigiComp} Grasegger, Georg RigiComp — A Mathematica package for computational rigidity of graphs, Zenodo, 2022 \href{https://alco.centre-mersenne.org/articles/10.5802/alco.390/#RigiComp}{Publisher reference}.
\bibitem[19]{LowerBounds} Grasegger, Georg; Koutschan, Christoph; Tsigaridas, Elias Lower bounds on the number of realizations of rigid graphs, Exp. Math., Volume 29 (2020) no. 2, pp. 125-136 \href{https://alco.centre-mersenne.org/articles/10.5802/alco.390/#LowerBounds}{Publisher reference}.
\bibitem[20]{GraverServatius} Graver, Jack; Servatius, Brigitte; Servatius, Herman Combinatorial rigidity, Graduate Studies in Mathematics, 2, American Mathematical Society, Providence, RI, 1993, x+172 pages \href{https://alco.centre-mersenne.org/articles/10.5802/alco.390/#GraverServatius}{Publisher reference}.
\bibitem[21]{Harris} Harris, Joe Algebraic geometry. A first course, Graduate Texts in Mathematics, 133, Springer-Verlag, New York, 1992, xx+328 pages \href{https://alco.centre-mersenne.org/articles/10.5802/alco.390/#Harris}{Publisher reference}.
\bibitem[22]{HopcroftTarjan} Hopcroft, John; Tarjan, Robert Algorithm 447: Efficient algorithms for graph manipulation, Communications of the ACM, Volume 16 (1973) no. 6, pp. 372-378 \href{https://alco.centre-mersenne.org/articles/10.5802/alco.390/#HopcroftTarjan}{Publisher reference}.
\bibitem[23]{Izm09} Izmestiev, Ivan Projective background of the infinitesimal rigidity of frameworks, Geom. Dedicata, Volume 140 (2009), pp. 183-203 \href{https://alco.centre-mersenne.org/articles/10.5802/alco.390/#Izm09}{Publisher reference}.
\bibitem[24]{JacksonOwen} Jackson, Bill; Owen, J. C. Equivalent realisations of a rigid graph, Discrete Appl. Math., Volume 256 (2019), pp. 42-58 \href{https://alco.centre-mersenne.org/articles/10.5802/alco.390/#JacksonOwen}{Publisher reference}.
\bibitem[25]{mumford} Mumford, David The red book of varieties and schemes, Lecture Notes in Mathematics, 1358, Springer-Verlag, Berlin, 1988, vi+309 pages \href{https://alco.centre-mersenne.org/articles/10.5802/alco.390/#mumford}{Publisher reference}.
\bibitem[26]{Pog73} Pogorelov, A. V. Extrinsic geometry of convex surfaces, Translations of Mathematical Monographs, Vol. 35, American Mathematical Society, Providence, RI, 1973, vi+669 pages (Translated from the Russian by Israel Program for Scientific Translations) \href{https://alco.centre-mersenne.org/articles/10.5802/alco.390/#Pog73}{Publisher reference}.
\bibitem[27]{roman} Roman, Steven Advanced linear algebra, Graduate Texts in Mathematics, 135, Springer, New York, 2008, xviii+522 pages \href{https://alco.centre-mersenne.org/articles/10.5802/alco.390/#roman}{Publisher reference}.
\bibitem[28]{Steffens2010} Steffens, Reinhard; Theobald, Thorsten Mixed volume techniques for embeddings of Laman graphs, Comput. Geom., Volume 43 (2010) no. 2, pp. 84-93 \href{https://alco.centre-mersenne.org/articles/10.5802/alco.390/#Steffens2010}{Publisher reference}.
\bibitem[29]{TayWhite85} Tay, Tiong-Seng; Whiteley, Walter Generating isostatic frameworks, Structural Topology (1985) no. 11, pp. 21-69 \href{https://alco.centre-mersenne.org/articles/10.5802/alco.390/#TayWhite85}{Publisher reference}.
\bibitem[30]{coning} Whiteley, Walter Cones, infinity and 1-story buildings, Structural Topology (1983) no. 8, pp. 53-70 \href{https://alco.centre-mersenne.org/articles/10.5802/alco.390/#coning}{Publisher reference}.
\end{thebibliography}
\end{document}
